% 01_intro.tex — owner: S (Round 5 reframe; supersedes the Round-4 version).
\section{Introduction}
\label{sec:intro}

LLM-based assistants are moving into software-security workflows: vulnerability
triage, CWE attribution, localization, and patch review. Function-level
vulnerability detection, however, remains far from solved on realistic
substrates: PrimeVul showed that results celebrated on legacy benchmarks
collapse under realistic evaluation, and paired metrics expose models that
cannot distinguish a vulnerable function from its patched
counterpart~\cite{ding2025primevul}. Systematic studies of LLMs echo this
pessimism once precision and realistic distributions are
considered~\cite{han2023howfar,sun2024llm4vuln}.

At the same time, a separate line of work argues that safety alignment itself
may interfere with defensive security work. \emph{Defensive Refusal Bias}
reports that cyber-defence requests containing security-sensitive vocabulary
are refused far more often than neutral paraphrases~\cite{campbell2026defensive};
over-refusal is measurable at scale~\cite{cui2025orbench,rottger2024xstest};
safeguards can even be weaponized into denial-of-service through their own
false positives~\cite{zhang2024safeguarddos}; and safety state has been linked
to vulnerability-analysis utility~\cite{li2026beyondrefusal}. Source code adds
a second channel of interference: comments, strings, and documentation are
model-facing text that an attacker can edit, and indirect prompt injection
(IPI) through such content is a demonstrated attack
surface~\cite{yi2023bipia,cheng2026codesentinel}. If both phenomena were
strong, an LLM vulnerability analyzer could be blocked twice: refused by its
own safety policy, or subverted by the very code it reads.

\textbf{A pre-registered reproduction gate.}
Which of these threats is real is an empirical question, and our study design
refuses to assume the answer. Before building any mitigation we registered a
reproduction gate (E0): if defensive refusal on vulnerability-analysis tasks
does not replicate under controlled prompt arms, the project must pivot ---
refusal becomes a secondary outcome and untrusted-context robustness becomes
primary~\cite{ding2025primevul,campbell2026defensive}. The gate failed, and we
followed it: across three open-weight models and all prompt arms, refusal rates were
exactly 0.000, while the same models refused 8--75\% of benign contrast
probes. Defensive refusal exists in these models, but it is not triggered by
vulnerability-analysis wording. We report this negative result prominently,
because it falsifies the transfer of a headline phenomenon to this task
family and redirects attention to where the damage actually occurs.

\textbf{What we measure instead.}
We present \emph{RefuseGuard}, a framework and pilot-scale study that measures
refusal, utility, and safety \emph{simultaneously} for LLM vulnerability
analysis under untrusted code context. On a semantics-preserving benchmark of
paired clean/contextualized PrimeVul functions (comments, docstrings, and
string carriers carrying security-charged but benign content, plus
instruction-like IPI content), availability metrics barely move:
usable-answer coverage stays at 0.98--1.00 and context-induced utility drop
(SIUD) is at most $+0.017$. The material effect is an \emph{integrity} one:
instruction-like context silently flips verdicts to ``benign'' on 74.2\% of
paired samples while the model still emits a well-formed answer.

\textbf{Escalating the attack --- and auditing the defence.}
A skeptic could attribute the null gate to weak stimuli: wording is easy to
see through. We therefore built C5, a content-anchored blocking attack in the
TabooRAG mould~\cite{li2026taboorag}: a short threat-intelligence advisory
synthesized from the analyzed function's own risky-sink fingerprint
(AST-extracted, label-blind) and embedded inside the artifact under analysis,
with audited anti-leakage (no label, CWE, or CVE leakage across all 400
renderings; balanced advisory strata within each label). The attack is
genuinely query-relevant where the naive context was not: 100/100 concrete
advisories name a real sink of the same function, versus 2/838 (0.2\%) naive
carriers. Pre-registered transfer gates over three models return
\textsc{not-supported} everywhere: still zero refusals. But the manipulation
is demonstrably not ignored --- it silently biases verdicts \emph{toward}
``vulnerable'' (Llama: 10--11 benign functions flipped per arm, never the
reverse; Granite: vulnerable recall $0.100\rightarrow0.733/0.767$). The
false-positive channel, not denial of service, is what a risk-context attack
actually opens at this scale.

Then we did to our own defence what an auditor should: P3, a provenance-based
safety wrapper that detects the advisory, labels it as untrusted data, and
reasserts the task intent, never refuses --- and never gets refused. It
simply \emph{corrupts the verdict}: Llama's vulnerable recall collapses from
$1.000$ to $0.367/0.333$ under the C5 arms ($39/60$ paired flips
vul$\to$benign, McNemar $p\leq3.8\times10^{-6}$), on audited, well-formed
JSON outputs. The defence, not the attack, is the strongest utility risk we
measured.

\textbf{Contributions.} Scoped honestly to pilot scale (hundreds to ~1.4
thousand new generations per study round, tens of samples per cell):
\begin{itemize}
  \item \textbf{C1 --- A query-relevant, anti-leakage blocking benchmark
  (C5)} for LLM vulnerability analysis: advisories synthesized from each
  function's own risky-sink fingerprint, label-blind by construction,
  stratum-balanced within each label, with pre-registered transfer gates and
  published manifests, seeds, and prompts (\S\ref{sec:method},
  \S\ref{sec:setup}).
  \item \textbf{C2 --- A three-model blocking-transfer study whose negative
  result is designed to be interpretable}: refusal does not transfer (all
  pre-registered gates \textsc{not-supported}, 0/3 models), and the
  anti-leakage audit plus an inert-evidence counterexample (the same
  manipulation moves verdicts) show the null is a property of the blocking
  mechanism at 2--3B scale --- a model-scale limitation we do not generalize
  from --- not a measurement artifact. The attack's measurable effect is a
  verdict bias toward ``vulnerable'' (\S\ref{sec:results}, RQ7).
  \item \textbf{C3 --- The defence-is-the-risk finding, with a flip-rate
  ledger}: the provenance wrapper P3 silently flips 39/60 vulnerable verdicts
  to benign on one model ($p\leq3.8\times10^{-6}$) while passing every
  availability and refusal metric; attributed at defence level only (P3 is a
  bundle of $\geq$5 components), and joined by the earlier result that
  ungated structured recovery manufactures unsafe compliance ($0.040\to
  1.000$) --- two independent ways a mitigation layer becomes the bug
  (\S\ref{sec:results}, RQ6--RQ7).
  \item \textbf{C4 --- Injection resistance with honest limits}: semantic
  isolation (P1) cuts aggregate instruction-injection success
  $0.742\rightarrow0.484$ ($p{=}0.0078$; vulnerable-only $p{=}0.5$) but does
  not improve accuracy --- isolation is not accuracy --- and a fine-tuned
  CodeBERT channel supplies refusal-independent coverage at realistic
  difficulty (F1 0.215) (\S\ref{sec:results}, RQ4--RQ5).
\end{itemize}

To our verified knowledge, this is the first systematic test of
query-relevant safety-blocking transfer to source-code vulnerability analysis
at open-model scale, and the first evidence of defence-induced verdict
corruption in this setting (\S\ref{sec:related} scopes both claims).

All numbers in this paper are traceable to released output artifacts; where a
run was partial or a measurement instrument failed validation, we say so and
keep the result in the threat model of the measurement itself
(\S\ref{sec:discussion}).
